\documentclass[a4paper,11pt]{article}
\usepackage[utf8]{inputenc}
\usepackage[T1]{fontenc}
\usepackage[english]{babel}
\usepackage[margin=22mm]{geometry}
\usepackage{lmodern}
\usepackage{amsmath,amssymb}
\usepackage{graphicx}
\usepackage{booktabs,tabularx,array}
\usepackage{microtype}
\usepackage{xcolor}
\usepackage{hyperref}
\hypersetup{hypertexnames=false,colorlinks=true,linkcolor=teal,urlcolor=teal,citecolor=teal,
  pdftitle={Geometric mode triplets and symmetry-compatible interactions},
\urlstyle{tt}
\setlength{\emergencystretch}{2em}
\newcommand{\C}{\mathbb C}
\newcommand{\I}{\mathrm I}
\newcommand{\norm}[1]{\left\lVert#1\right\rVert}
\newenvironment{drafttable}{\begin{center}\small\renewcommand{\arraystretch}{1.2}}{\end{center}}

\begin{document}
\begin{flushleft}
Not peer reviewed \textperiodcentered\ Authorship and publication version remain to be determined\\
English translation of version 0.5}

\vspace{1em}
{\LARGE\bfseries Geometric mode triplets and symmetry-compatible interactions\par}
\vspace{0.7em}
{\bfseries Analytical constructions, numerical cross-checks, and limits of an SU(3) interpretation}
\end{flushleft}

\section*{Abstract}
We investigate the conditions under which finite geometric oscillator models possess and retain a three-dimensional complex mode sector with continuous internal symmetry. The Laplacians of the tetrahedral and cube graphs contain degenerate triplets. Their quadratic amplitude dynamics has U(3) symmetry with an SU(3) subgroup. An ordinary local quartic nonlinearity, however, destroys this symmetry already within the triplet and generally couples in additional modes. We give conditions for symmetry-compatible exchange between two triplets, construct a conservative total-density mediator, and test geometric countermodels. None of the 54 individual distance kernels investigated passed the predefined isotropy gate. Identical matched connections and a symmetric sum over six spatial directions, by contrast, provide controlled isotropic constructions. The results concern specified reduced models; a derivation from spatial Q-ball fields, local gauge dynamics, and a QCD interpretation remain outstanding.

\section*{1. Question and scope of the claims}
The starting point is the question of whether four or eight coupled constituents can produce three equivalent internal states and suitable interactions. We distinguish the existence of a degenerate linear mode sector from its nonlinear preservation and from a local gauge theory. We study dimensionless complex graph amplitudes. The geometry is prescribed; the constituents do not move together on their own to form a bound body.

The mathematical tools---graph spectra, unitary invariance, projection, and completing the square---are established. No claim to novelty or priority is made. The contribution of this working draft consists in their concrete combination, the counterchecks, and the identification of the missing steps in the derivation.

\begin{drafttable}
\begin{tabularx}{\linewidth}{@{}lXX@{}}
\toprule
Quantity & Meaning & Not to be equated with\\\midrule
$N=4$ or $8$ & Number of graph sites & Colors or gluons\\
$a\in\C^3$ & Three complex mode amplitudes & Already established quark states\\
$l,m$ & Angular momentum or winding index in other project models & Dimension of this triplet\\
$A=1,\ldots,8$ & Generator index of SU(3) & Eight spatial vertices\\\bottomrule
\end{tabularx}
\end{drafttable}

\section*{2. A geometrically generated linear triplet}
For $z\in\C^N$, we choose canonical complex amplitudes and the graph Laplacian $L$:
\begin{equation}
H_2=z^\dagger Lz,\qquad i\dot z=Lz. \tag{1}
\end{equation}
This canonical structure is a model assumption. It is not obtained here from a relativistic field theory. For the complete tetrahedral graph $K_4$ and the edge graph of the cube $Q_3$, the following hold exactly:
\begin{equation}
\begin{aligned}
\operatorname{spec}(L_{K_4})&=\{0,4,4,4\},\\
\operatorname{spec}(L_{Q_3})&=\{0,2,2,2,4,4,4,6\}.
\end{aligned}\tag{2}
\end{equation}
Let $X$ be the matrix of vertex coordinates with entries $\pm1$: for the tetrahedron, the rows are $(1,1,1)$, $(1,-1,-1)$, $(-1,1,-1)$, and $(-1,-1,1)$; for the cube, they are all eight sign combinations. Then $U=X/\sqrt N$ has orthonormal columns. For $z=Ua$, we obtain $H_2=\lambda a^\dagger a$, with $\lambda=4$ and $2$, respectively. Every transformation $a\mapsto Wa$, with $W\in\mathrm{U}(3)$, preserves this functional.

\textbf{Consequence.} The degenerate complex quadratic sector has an additional continuous internal symmetry. The spatial symmetry group of the finite body remains discrete. More precisely, $\mathrm{U}(3)\cong[\mathrm{SU}(3)\times\mathrm{U}(1)]/\mathbb Z_3$. The common phase has therefore not been removed.

Increasing a single edge weight from $1$ to $1.01$ lifts the degeneracy: the tetrahedral triplet becomes $\{4,4,4.02\}$, and the lower cube triplet becomes $\{2,2,2.0049688749\}$. Moreover, the uniform mode with eigenvalue zero lies \emph{below} the triplet under consideration. A low-energy argument alone therefore does not justify using it in isolation.

\section*{3. Nonlinearity as a countercheck}
With $S=a^\dagger a$, projecting a local quartic energy gives exactly
\begin{equation}
\sum_i |(Ua)_i|^4=\frac{2S^2+|a\cdot a|^2-2\sum_\alpha|a_\alpha|^4}{N}.\tag{3}
\end{equation}
The identity follows by expansion and from the second and fourth sign moments of the vertex coordinates. At $S=1$, its value is $1/N$ for both $a=(1,0,0)$ and $(1,i,0)/\sqrt2$, but $7/(3N)$ for $(1,1,1)/\sqrt3$. States of equal norm thus have different energies: the linear U(3) symmetry does not survive this local nonlinearity.

In addition, $|Ua|^2Ua$ generally has a component outside the triplet. Projected onto the normalized external mode, its amplitude is
\begin{equation}
b_{\mathrm{out}}=\frac{2}{N}(a_xa_y\bar a_z+a_xa_z\bar a_y+a_ya_z\bar a_x).\tag{4}
\end{equation}
For the tetrahedron, this is the uniform mode; for the cube, it is the $xyz$ parity mode. Particularly symmetric one- or two-component initial states can hide this effect.

Numerically, the full equation $i\dot z=Lz+0.5|z|^2z$, without projection back into the triplet, was integrated to $T=20$. The initial vector in the triplet was proportional to $(1,0.7+0.3i,-0.2+0.6i)$. DOP853 with tolerances $10^{-9}$ and $10^{-11}$, a maximum step of $0.1$, and $401$ output points gave the following maximum relative leakage over the $401$ output points at the finer accuracy level:

\begin{drafttable}
\begin{tabular}{@{}lll@{}}\toprule
Initial amplitude & Tetrahedron & Cube\\\midrule
$0.1$ & $3.861\times10^{-8}$ & $9.512\times10^{-9}$\\
$0.3$ & $3.918\times10^{-6}$ & $8.502\times10^{-7}$\\
$1$ & $8.910\times10^{-4}$ & $2.021\times10^{-4}$\\\bottomrule
\end{tabular}
\end{drafttable}
The leakage is $\norm{(I-UU^\dagger)z}^2/\norm{z(0)}^2$. The norm and energy gates of $10^{-7}$ and the final-state comparison gate of $10^{-6}$ were passed. Small leakage does not guarantee symmetry: equation (3) already breaks the symmetry within the sector.

\section*{4. Which exchange interaction is compatible?}
Let two graphs be coupled by nonnegative weights $W_{ij}$:
\begin{equation}
\begin{aligned}
H_{\mathrm{int}}&=\sum_{ij}W_{ij}|z_i-z'_j|^2\\
&=a^\dagger Aa+b^\dagger Bb-a^\dagger Cb-b^\dagger C^\dagger a,\\
A&=U^\dagger D_AU,\qquad B=V^\dagger D_BV,\qquad C=U^\dagger WV.
\end{aligned}\tag{5}
\end{equation}
The second line of (5) applies to $z=Ua$ and $z'=Vb$. The matrices $D_A$ and $D_B$ contain the row and column sums of $W$, respectively. In a common internal basis, this expression is invariant under every common SU(3) transformation if and only if $A=\alpha I$, $B=\beta I$, and $C=cI$. This follows from the commutant of the irreducible fundamental representation. A relative unitary choice of basis can make $C$ a scalar matrix only if its three singular values are equal; it does not repair anisotropic self-blocks.

Furthermore, the sector must be preserved in the full dynamics. With $P_U=UU^\dagger$ and $P_V=VV^\dagger$, the following are additionally required:
\begin{equation}
\begin{aligned}
(I-P_U)D_AU&=0,\qquad &(I-P_V)D_BV&=0,\\
(I-P_U)WV&=0, &(I-P_V)W^{\mathsf T}U&=0.
\end{aligned}\tag{6}
\end{equation}
\textbf{Matched connections.} For identical bases and $W=JI$, one obtains exactly $H_{\mathrm{ex}}=J\norm{a-b}^2$ without leakage. This contact law is explicitly imposed, rather than derived from free field superposition. In contrast, $W=w\mathbf1\mathbf1^{\mathsf T}$ gives $C=0$ for the nonuniform triplets: a uniform connection of all sites shifts frequencies but generates no triplet exchange.

\subsection*{4.1 Distance kernels and six spatial directions}
As a countermodel, we chose $W_{ij}=\exp[-|r_i-r'_j|^2/(2\ell^2)]$. A vertex radius of $1$, separation $d\in\{3,4,6\}$, range $\ell\in\{1,2,4\}$, and three orientations (identity, $R_z(\pi/4)$, $R_y(\pi/2)$) give $54$ cases across both geometries. No case passed the relative isotropy gate of $10^{-10}$ with $\norm C>10^{-8}$. The smallest cross-block anisotropy value was $0.29739$. This is a finite search result, not an exclusion of all distance laws.

The sum of identical kernels for the six displacements $\pm4e_x$, $\pm4e_y$, and $\pm4e_z$ at $\ell=2$, however, gave $C=-0.05845231382I$ (tetrahedron) and $C=-0.05960274201I$ (cube), with isotropic self-blocks and leakage residuals below $1.1\times10^{-15}$. Negative cross-blocks do not contradict the positivity of the original spring energy.

\begin{quote}
\textbf{Limit of validity:} An aggregated coupling kernel was tested. Six independently moving neighbors constitute a larger dynamical system. Their coherent combination must be examined there with its own normalization (a factor of $\sqrt6$ for six equal amplitudes) and with all orthogonal neighbor modes. The static sum does not prove its dynamical preservation.
\end{quote}

\section*{5. A conservative model with frequency feedback}
An explicitly U(3)-compatible nonlinearity depends only on the total density $S$. For two triplets, each with one real mediator, we define
\begin{equation}
\begin{aligned}
H={}&\omega(S_A+S_B)+J\norm{a-b}^2\\
&+\sum_{X=A,B}\left[\frac{\lambda S_X^2}{2}+\frac{P_X^2}{2}
+\frac{Ky_X^2}{2}+hy_XS_X\right].
\end{aligned}\tag{7}
\end{equation}
Completing the square gives, for each constituent,
\begin{equation}
\frac{K}{2}(y+hS/K)^2+\frac{\lambda-h^2/K}{2}S^2+\frac{P^2}{2}.\tag{8}
\end{equation}
For positive $\omega,J,K$ and $\lambda>h^2/K$, the energy is bounded below. With $\omega=K=1$, $\lambda=0.3$, and $h=0.4$, the residual quartic coefficient is $0.14$. Static minimization over the mediator is exact; using it as a dynamical elimination would additionally require a separation of timescales.
\begin{equation}
\begin{aligned}
i\dot a&=(\omega+\lambda S_A+hy_A)a+J(a-b),\\
i\dot b&=(\omega+\lambda S_B+hy_B)b+J(b-a),\\
\dot y_X&=P_X,\qquad \dot P_X=-Ky_X-hS_X.
\end{aligned}\tag{9}
\end{equation}
Direct differentiation shows that $M=aa^\dagger+bb^\dagger$ is conserved. Thus the total norm $\operatorname{tr}M$ and the eight quantities $Q_A=\operatorname{tr}[(\lambda_A/2)M]$ are conserved. Here $\lambda_A$ denotes the Gell-Mann matrices, not the quartic coefficient. These are Noether quantities of the model, not already identified electric or QCD color charges.

For an isolated triplet, the mediator alone changes only its common phase velocity. Population exchange between the constituents becomes possible only for $J\ne0$. For $J\in\{0,0.15\}$ and $h\in\{0,0.4\}$, four arms were computed to $T=80$, each at two accuracy levels. The initial data before common normalization were $a=(1,0,0)$ and $b=(0,0.7,0.3i)$. Both vectors are divided by $\sqrt{1.58}$. In addition, $y=(0,0)$ and $P=(0.1,-0.1)$.

\begin{drafttable}
\begin{tabular}{@{}lll@{}}\toprule
$J$ & $h$ & Range of $|a_1|^2$ up to $T=80$\\\midrule
$0$ & $0$ or $0.4$ & Constant at $0.6329113924$\\
$0.15$ & $0$ & $0.0000006121$ to $0.6329113924$\\
$0.15$ & $0.4$ & $0.0000365731$ to $0.6329113924$\\\bottomrule
\end{tabular}
\end{drafttable}
In the last arm, the maximum energy drift at the finer accuracy level was $1.17\times10^{-12}$, and the charge-matrix drift was $6.12\times10^{-13}$. The stated population ranges refer to $801$ output points; continuous-time extrema were not certified. The tolerance comparisons passed. This closed dynamics contains neither radiative loss nor free spatial particle motion. Long-term spatial binding or a bound state in the continuum (BIC) does not follow from it.

\subsection*{5.1 New countercheck: six independent neighbors}
Following the first draft, the missing linear network test was performed. A central cluster and six independent neighboring clusters at $\pm4$ along the coordinate axes each have their own internal graph Laplacian. Only the center and neighbors are coupled. The resulting Hamiltonian matrices have dimensions $28$ and $56$ for the two geometries. Identical matched connections with $J=1$ are compared with the distance kernel from Section 4.1.

Two subspaces must be distinguished: the $21$ independent local triplet amplitudes, and the six-dimensional space consisting of the central triplet plus an equal neighbor amplitude. In the second space, the canonically normalized embedding is $z_0=Ua$, $z_s=Ub/\sqrt6$. For matched connections, this space is exactly invariant; its Hamiltonian is
\begin{equation}
H_{\mathrm{collective}}=
\begin{pmatrix}\lambda+6J&-\sqrt6J\\-\sqrt6J&\lambda+J\end{pmatrix}\otimes I_3.\tag{10}
\end{equation}
For the distance kernel, too, the merely projected $6\times6$ matrix is isotropic (numerical residual below $1.3\times10^{-15}$). Its block coupling to the exterior, however, does not vanish. The Frobenius norm of this block is $0.88009$ for the tetrahedron and $1.79607$ for the cube. The full time evolution therefore leaves the projected space.

\begin{drafttable}
\begin{tabularx}{\linewidth}{@{}XXX@{}}\toprule
Geometry / initial state & Outside local triplets & Outside coherent combination\\\midrule
Tetrahedron / central & $1.382\%$ & $2.802\%$\\
Cube / central & $4.676\%$ & $6.491\%$\\
Tetrahedron / neighbor sum & $0.605\%$ & $99.749\%$\\
Cube / neighbor sum & $23.242\%$ & $96.333\%$\\\bottomrule
\end{tabularx}
\end{drafttable}
The entries are maximum norm fractions over $801$ output points up to $T=80$. The matched connection showed only residuals below $6.3\times10^{-29}$ for all initial states. Time evolution was evaluated by eigendecomposition of the full linear Hamiltonian; norm and energy drift remained below $3\times10^{-14}$. The new run required $0.047$ CPU-s on TS440.

The mechanism is already apparent analytically: the uniform mode of an individual neighbor receives the amplitude $u_0^{\mathsf T}W_s^{\mathsf T}Ua$. These contributions can cancel in an aggregated sum, but in the full network they lie in different neighbor spaces. Their squared norms do not cancel. This disproves the transfer of static sum isotropy to the chosen dynamical reduction. It does not exclude other globally distributed degenerate eigenmodes of the entire network.

\subsection*{5.2 Field mechanism: radial exchange in concentric shells}
An additional construction yields equal channel couplings from a single scalar field rather than from three matched ports. We choose the linear real Klein--Gordon action with a prescribed radial potential:
\begin{equation}
\begin{aligned}
L&=\frac12\int[\dot\phi^2-|\nabla\phi|^2-(10+V(r))\phi^2]\,d^3x,\\
V(r)&=-4\exp\left[-\left(\frac{r-6}{0.6}\right)^2\right]
-4\exp\left[-\left(\frac{r-6-d}{0.6}\right)^2\right].
\end{aligned}\tag{11}
\end{equation}
This is a new linear test model with external shells, not a self-consistent Q-ball solution. With $\phi=\sum_i u_i(r,t)Y_i(\Omega)/r$ and the real normalized dipole harmonics $Y_i=\sqrt{3/(4\pi)}\,n_i$, angular integration gives the same radial operator $D_1=-\partial_r^2+2/r^2+10+V(r)$. The radial measure is $dr$, with $\int u_nu_k\,dr=\delta_{nk}$ for the eigenfunctions. Regularity at the origin and decay on the half-line are implemented numerically by a regular origin and a Dirichlet termination of the box. The identical angular treatment follows from spherical symmetry and orthogonality.

Two exact radial eigenfunctions with positive eigenvalues $e_1,e_2$ generate an invariant linear sector. Canonical oscillator amplitudes $c_n=(\sqrt{\omega_n}q_n+ip_n/\sqrt{\omega_n})/\sqrt2$, with $\omega_n=\sqrt{e_n}$, give $H=\sum_n\omega_n|c_n|^2$ without a rotating-wave approximation. A common radial basis rotation $O$, determined by diagonalizing the radius operator projected into this sector, gives
\begin{equation}
\begin{aligned}
H&=\varepsilon_Aa^\dagger a+\varepsilon_Bb^\dagger b-J(a^\dagger b+b^\dagger a),\\
K&=O^{\mathsf T}\operatorname{diag}(\sqrt{e_1},\sqrt{e_2})O,\qquad J=-K_{12}.
\end{aligned}\tag{12}
\end{equation}
The exchange strength is thus obtained from the radial field calculation and is the same for all three angular components. General U(3) transformations of the canonical amplitudes are not simply spatial rotations; they can mix field coordinates and momenta.

\begin{drafttable}
\begin{tabularx}{\linewidth}{@{}l l X l@{}}\toprule
Shell separation $d$ & $|J|$ & Maximum canonical transfer & First maximum $t$\\\midrule
$2$ & $0.0637380$ & $99.909\%$ & $24.633$\\
$3$ & $0.0172399$ & $97.435\%$ & $89.938$\\
$4$ & $0.00471585$ & $68.151\%$ & $274.976$\\\bottomrule
\end{tabularx}
\end{drafttable}
The values come from the grid with $R=24$ and $h=0.025$. Halving the grid spacing at $R=20$ changed $|J|$ by less than $3\times10^{-4}$ in relative terms; enlarging the box from $20$ to $24$ changed it by less than $7\times10^{-11}$. The norm, eigen-residuals, analytical two-mode time evolution, and predefined convergence gates passed. Nine radial cases and the angular controls required $0.0183$ CPU-s on TS440.

\textbf{Interpretation of the percentages:} They refer to $|b|^2$ for an initially normalized inner canonical combination. This is neither the fraction of spatial field energy in the outer shell nor a particle quantity. Physical field coordinates carry frequency-dependent factors of $1/\sqrt{\omega_n}$. Transfer is determined analytically from the two-mode matrix; the numerical eigenvalues are not interval-certified.

Nonlinearity remains the limitation. For a complex spatial dipole amplitude $\eta=a\cdot Y$, the identity $\int|\eta|^4\,d\Omega=3[2(a^\dagger a)^2+|a\cdot a|^2]/(20\pi)$ holds exactly. A linear and a circular normalized amplitude therefore differ by a factor of $1.5$. This angular test disproves automatic U(3) invariance of a local complex quartic term. It is not an already derived nonlinear normal form of the real field in (11). An additional $P_2$ angular perturbation also splits the projected operator itself and couples in $l=3$.

Kaminaga's concentric $\delta$-shells provide a mathematical parallel result for the $m$-independent partial-wave structure [3, Lemma 9 and Remark 4]. His asymptotic tunneling theorem, however, concerns matched limiting levels at $l=0$ and large separation [3, Theorem 17]. It establishes neither our smooth $l=1$ numerical values nor exact complete transfer at finite separation; the latter follows here from our own two-mode algebra. Technical review notes on the printed proof in \S\S6--7 are documented in the source audit; its error orders are not adopted as a certificate. For the original Q-balls, the coupled $l=1$ problem remains decisive. Panin and Smolyakov identify translations in a different two-field model [4, Eq.~(22)]; their stability study does not prove the absence or presence of positive internal modes in our model.

\begin{quote}
\textbf{Result and remaining gap:} Equal linear exchange can be derived from a common spatial field geometry. The autonomous emergence of this geometry and preservation of the symmetry under the actual nonlinear field couplings remain open. In our multicomponent overall formula, internal U(3) at $N=3$ and $g=0$ would already be assumed, rather than explained by spatial geometry; the existing $g$ term generally breaks it.
\end{quote}

\subsection*{5.3 What the 68\% means: detuning and spatial energy}
The maximum transfer at shell separation $d=4$ is not a special universal number. For the canonical two-mode matrix in (12), with $\delta=\varepsilon_A-\varepsilon_B$, we have
\begin{equation}
P_B(t)=\frac{4J^2}{\delta^2+4J^2}\sin^2\left[\frac{\sqrt{\delta^2+4J^2}\,t}{2}\right].\tag{13}
\end{equation}
$J$ controls exchange, and $\delta$ controls frequency detuning. As the separation increases, $J$ decreases, so a small residual detuning limits transfer more strongly. The untransferred fraction is neither missing energy nor a separate dark sector. Cosmological quantities do not occur in this model.

A new control changed only the depth $A$ of the outer well, keeping the radii fixed at $6$ and $10$. A predefined grid $A\in[3.9,4.1]$ and a root search for $\delta(A)$ restricted to this interval gave $A=3.9375147759$ on the finest grid. Compared with $A=4$, this is a reduction of about $1.56\%$. Canonical transfer then analytically reaches its full value at $\delta=0$, provided $J\ne0$. Numerically, $|\delta|<5\times10^{-14}$. This is an expected control of frequency matching, not a new constant of nature.

In addition, the real field was reconstructed and its local energy for $r>8$ was calculated. The radial physical energy density is $\tfrac12[\dot u^2+(u'-u/r)^2+2u^2/r^2+(10+V)u^2]$. Relative to the reduced radial functional, the inner partial energy receives the boundary term $-u(8)^2/16$, and the outer energy receives the opposite term. The discrete total energy thus remains exactly consistent with the field operator used.

\begin{drafttable}
\begin{tabularx}{\linewidth}{@{}lXXl@{}}\toprule
Case, $h=0.0125$ & Max. mode transfer & Outer field energy at mode maximum & Time\\\midrule
$A=4$ & $68.153\%$ & $67.934\%$ & $274.965$\\
$A=3.9375147759$ & $100\%$ at frequency matching & $99.635\%$ & $328.715$\\\bottomrule
\end{tabularx}
\end{drafttable}
Initially, about $0.54\%$ of the field energy is already outside the cut. The second table row therefore corresponds to a net increase of about $99.09$ percentage points of the total energy in the outer region. The energy is reported at the \emph{canonical} maximum; fast carrier phases can shift the spatial energy maximum slightly. The canonical excitation returns later. Incomplete transfer does not imply irreversible dissipation in this conservative model.

\begin{figure}[htbp]
\centering
% Source figure alternative text: Detuning changes the maximum transfer;
% time curves show the reversible exchange.
\includegraphics[width=\linewidth]{transfer-detuning.pdf}
\caption{Outer well depth and exchange. Lines on the right: canonical mode populations; crosses: three explicitly evaluated spatial energy fractions. No interpolation of fast energy oscillations. Units are dimensionless.}
\end{figure}

Computing location: TS440, $0.0774$ CPU-s. Three grid spacings, $h=0.05$, $0.025$, and $0.0125$ at $R=24$, separate convergence gates, and numerical conservation checks passed. With grid halving, the additional direct midpoint quadrature approaches the regional energy accounting at approximately quadratic order. No separate new box-size test was performed for the matched depth. Source code, parameters, and data: \path{coordination/resonance-20260930/shell-detuning/}.

\subsection*{5.4 Test on the original Q-ball and a conditional chain of proof}
The external shell mechanism was subsequently checked for its prerequisites in the actual single-field model $U(S)=S-S^2+S^3/2$. We used our own previously stored profile at $\omega^2=0.7$ ($Q\approx473.413$), not the earlier Higgs-portal Q250 candidate. For the reduced radial $l=1$ perturbation, the operators are
\begin{equation}
\begin{aligned}
A&=-\partial_r^2+2/r^2+0.3-6f^2+7.5f^4,\\
B&=-\partial_r^2+2/r^2+0.3-2f^2+1.5f^4,\\
M(\rho)&=\begin{pmatrix}A-\rho^2&-2\omega\rho\\-2\omega\rho&B-\rho^2\end{pmatrix}.
\end{aligned}\tag{14}
\end{equation}
On the exact half-line, the translation mode satisfies $A(rf')=0$. It is a displacement of the entire object, not an additional positive internal oscillator. Both free sidebands remain closed simultaneously only for $0<\rho<1-\omega\approx0.16333997$.

\textbf{Our own conditional analytical statement.} If $B\ge bI$ and $0\le s=\rho^2<b$, the lower component can be eliminated exactly. The Schur operator and its derivative are
\begin{equation}
\begin{aligned}
F(s)&=A-sI-4\omega^2s(B-sI)^{-1},\\
F'(s)&=-I-4\omega^2B(B-sI)^{-2}\le-I.
\end{aligned}\tag{15}
\end{equation}
No commutativity of $A$ and $B$ is required. Congruence by Schur elimination preserves the negative inertia because $B-sI$ is positive. All ordered eigenvalues of $F$ therefore decrease strictly with $s$. Under established initial-index, gap, and endpoint conditions, the increase in the negative index counts the additional positive-frequency zeros. For a singular right endpoint, its kernel dimension must also be counted. Simply projecting the translation out orthogonally would generally be wrong; the full coupled inertia avoids this error.

\begin{drafttable}
\begin{tabularx}{\linewidth}{@{}lXXX@{}}\toprule
Box / step & Translation eigenvalue of $A$ & $\lambda_{\min}(B)$ & Second eigenvalue of $M$ at test endpoint\\\midrule
$R=24$ / $h=0.08$ & $-1.5391\times10^{-4}$ & $0.2131403$ & $0.0398112$\\
$R=32$ / $h=0.04$ & $-3.8457\times10^{-5}$ & $0.2131733$ & $0.0211452$\\
$R=40$ / $h=0.02$ & $-9.6131\times10^{-6}$ & $0.2131817$ & $0.0132527$\\\bottomrule
\end{tabularx}
\end{drafttable}
The test endpoint is $\rho=1-\omega-10^{-4}$. In all three discrete box problems, the negative index remains at one, and the endpoint is numerically nonsingular. No additional bound dipole mode was identified. The small negative translation eigenvalue shrinks approximately quadratically with $h$ and is not interpreted as a physical instability. The free negative control passes. Computing location: TS440, $2.558$ CPU-s; software testing beforehand on .69.

\begin{quote}
\textbf{No claim of exclusion in the continuum:} The finite box shifts weakly bound states. At the test endpoint, the slow asymptotic decay length is about $70.7$, exceeding the largest box radius of $40$. The second positive eigenvalue also decreases as the box grows. Very extended weak binding, the final frequency layer of $10^{-4}$, the exact threshold, and open resonances therefore remain undecided. The inertia counts are not interval-certified. The shell test remains a working external control model, but the two positive dipole branches it requires have not yet been demonstrated in the original Q-ball.
\end{quote}

The next proof task therefore concerns controlled treatment of the exterior, threshold behavior, and validated spectral counting. Our Schur derivation and its prerequisites are given in the accompanying source audit; numerical artifacts are located in \path{coordination/resonance-20260930/qball-l1-gap/}.

\section*{6. Relation to field theory and the literature}
Hammer studies singly excited states of identical, fixed atoms with parallel dipoles on a regular ring, optionally with a central atom. The cyclic group $\mathbb Z_N$ organizes the states; an effective non-Hermitian operator describes shifts and decay in the one-photon sector under consideration within the rotating-wave approximation [1, \S\S II--IV]. Here, subradiant means suppressed decay, not automatically an exact BIC [1, \S VI]. This is a limited analogy for mode organization, not a derivation of our conservative graph model or its SU(3) interpretation.

In QCD, quark fields carry a fundamental SU(3) color index; eight gluon components form the gauge connection. The field strength contains the non-Abelian structure-constant term from which gluon self-interactions follow [2, \S9.1, Eqs.~(9.1)--(9.2)]. Our model so far has only a common global symmetry. An algebraic distinction of our own is the following: for globally uniquely chosen complete basis matrices $G_i$ and links $L_{ij}=G_i^\dagger G_j$ formed exclusively from them, every closed link product telescopes to $I$. Moving projected subspaces, additional connections, or nontrivial transition functions are not covered by this statement. This basis-link calculation is our derivation, not a PDG theorem checked here.

\section*{7. What is missing for the further derivation}
\begin{enumerate}
\item \textbf{A common microscopic action.} The initial fields, couplings, boundary conditions, and stationary many-body configurations must be specified. The graph model must follow from this action rather than being postulated alongside it.
\item \textbf{A canonical field-to-mode mapping.} Linearization and symplectic normalization must yield an actual complex triplet. For rotating Q-balls, the coupled frequency problem must be treated; a spatial Hessian block alone is insufficient. The lower singlet mode must not be silently omitted.
\item \textbf{A controlled nonlinear reduction.} The full quartic coefficient tensor, couplings to excluded modes, and errors over time must be derived. Either the total-density form emerges, or the symmetry breaking must be quantified. Equation (3) is the mandatory countercheck.
\item \textbf{A microscopic exchange mechanism.} The matrices $A,B,C$ and the four leakage blocks must be calculated from field overlaps or a justified mediator. Matched ports or six spatial directions must not be chosen merely to obtain the desired result.
\item \textbf{The complete neighbor network.} The central cluster plus six independent neighbors must be calculated including their relative phases, complementary modes, and perturbations. The linear star was tested in Section 5.1: the distance kernel fails invariance; matched connections pass. Nonlinear networks and derived field couplings remain open.
\item \textbf{Local rather than merely global symmetry.} Independently variable connection degrees of freedom, a derived action, and its Gauss constraints would be required. The theory must allow nontrivial curvature without requiring every state to be curved. The additional U(1) sector must be explicitly accounted for.
\item \textbf{Only then, a QCD test.} Quantization, statistics, a relativistic limit, and suitable gauge observables are missing. Eight generators or three amplitudes replace neither gluon dynamics nor confinement nor asymptotic freedom.
\end{enumerate}
\textbf{Priority:} Test steps 2--5 first. Failure of the nonlinear triplet reduction or of network invariance would already be a clear finding against the respective mechanism. Adding further gauge fields before this bridge holds would merely define a different model.

\section*{8. Reproducibility and limitations}
The reported small physics calculations ran on TS440 with one CPU core and the shared laboratory limit. The triplet suite required $0.843$ CPU-s, and the interaction suite $4.823$ CPU-s; the supplementary leakage test is logged separately. These are floating-point calculations with convergence and conservation checks, not rigorous interval certificates. No historical physics runs were repeated for this text.

Source code, parameters, results, and the original input/output hashes are in the following project directories. The source-file hashes associated with this draft are additionally recorded in \path{SOURCE-SHA256SUMS.txt}:
\begin{itemize}
\item \path{coordination/resonance-20260930/su3-triplets/}: \path{PLAN.txt}, \path{run.py}, \path{RESULT.json}, \path{ERGEBNIS.txt}.
\item \path{coordination/resonance-20260930/su3-interactions/}: \path{PLAN.txt}, \path{run.py}, \path{RESULT.json}, \path{leakage.py}, \path{LEAKAGE.json}, \path{ERGEBNIS.txt}.
\end{itemize}
The analytical statements are conditional consequences of the specified models. Numerical control residuals are not experimental error bars. No model of natural particles, local SU(3) gauge theory, QCD, or Millennium proof is claimed.

\section*{References}
\begin{enumerate}
\renewcommand{\labelenumi}{[\theenumi]}
\item H. Hammer, \emph{Collective excitations in circular atomic configurations, and single-photon traps}, Physical Review A \textbf{70}, 023803 (2004). \href{https://arxiv.org/abs/quant-ph/0402065}{arXiv:quant-ph/0402065}. Checked: v5, \S\S II--IV and VI, especially Eqs.~(14)--(24) and (38)--(39).
\item J. Huston, K. Rabbertz, G. Zanderighi, \emph{Quantum Chromodynamics}, Particle Data Group, revision August 2025, Section 9.1. \href{https://pdg.lbl.gov/2025/reviews/rpp2025-rev-qcd.pdf}{PDG QCD Review}. Checked: \S9.1, pp.~1--2, Eqs.~(9.1)--(9.2).
\item M. Kaminaga, \emph{Schr\"odinger operators with concentric $\delta$--shell interactions}, arXiv:2602.09376v1 (2026). \href{https://arxiv.org/html/2602.09376v1}{Primary text}. Checked: \S\S2--7, Lemma 9, Remark 4, Theorem 17, and Appendix B; documented technical review notes on \S\S6--7 in \path{SOURCE-3.txt}.
\item A. G. Panin, M. N. Smolyakov, \emph{Classical behaviour of Q-balls in the Wick--Cutkosky model}, Eur. Phys. J. C \textbf{79}, 150 (2019). \href{https://link.springer.com/article/10.1140/epjc/s10052-019-6638-2}{Primary text}. Checked: \S\S1, 3.1, Eqs.~(12), (18)--(22), and (29)--(31); a different two-field potential, with no transfer of its spectral finding.
\end{enumerate}

\bigskip\hrule\medskip
{\small Working status: analytical model statements and limited numerical cross-checks. A full literature novelty review and external peer review remain outstanding.}
\end{document}
